\documentclass[12pt]{article}
\pagestyle{empty}
\topmargin=0mm
\oddsidemargin=0mm
\textwidth170mm
\textheight250mm

\begin{document}

\begin{center}
{\bf APPLICATION OF COMPUTER ALGEBRA SYSTEMS TO CONSTRUCT FORMULAE OF THE
COLLOCATION AND LEAST SQUARES METHOD}\\
\vspace {4mm}

L.G.~Semin, V.P.~Shapeev\\
{\it \small ITAM SB RAS, 630090, Novosibirsk, Russia\\
e-mail: shapeew@itam.nsc.ru}
\end{center}
\vspace {6mm}

Huge analytical work is needed when constructing the formulae of the
collocation and least squares (CLS) method for solving the Navier\,---\,Stokes
equations numerically.  To great extent, such work can be carried out by
computer algebra systems (CAS) \cite{bk:val, bk:sha}.  To construct the
formulae one needs to perform analytically the follows: summation of huge
expressions in functionals, their differentiation, solving a linear algebraic
systems in order to obtain coefficients of approximate solution in polynomial
basis.  CAS can successfully cope with this work.  Computer derivation of the
formulae of the method also gives additional advantages.  Optimization with
respect to rate of computation by obtained formulae, analytical and numerical
verification of these formulae by using the same or another CAS, conversion of
constructed formulae to arithmetic operators of FORTRAN and C programming
languages can be made.  All this greatly facilitate the final result obtaining
and reduce the time consumption for it.  Different modifications of the method
are easily available using computer programs in CAS.  Modifications related to
replacement of basis, collocation nodes, matching conditions  and so on, can
be made just by changing functions involved in the program.  Construction of
the formulae of CLS method for Navier\,---\,Stokes equations was first
performed with the help of the computer program of analytic computations in
CAS REDUCE, and then in MAPLE \cite{bk:sha, bk:semin, bk:hungary}.

The CLS method created was verified by solving the different problems
numerically.  Comparison between obtained and known results was made, good
potentialities of the method were shown.  This work was partially supported
by the Russian Foundation for Basic Research (grants 99--01--00515 and
00--01--00370).

\begin{thebibliography}{}
\bibitem{bk:val}
A.N. Valiullin, V.P. Shapeev et al. 
Symbolic manipulations in the
methods of mathematical physics //
Symposium mathematics for computer science.
Paris, March 16--18, 1982, P.~431-438.

\bibitem{bk:sha}
V.P. Shapeev.
Application of computer algebra in numerical methods //
Proceedings of 31-th regional youth conference "Problems of theoretical
and applied mechanics" --- Ekaterinburg, 2000. --- Pp.~80--84
(in Russian)

\bibitem{bk:semin}
  L.G.~Semin, V.P.~Shapeev.
  Collocation and least squares method for Navier\,---\,Stokes equations //
  Computational Technologies. --- 1998. --- Vol.~3, No.~3. --- Pp.~72--84.
  (in Russian)

\bibitem{bk:hungary}
  L.G.~Semin, V.P.~Shapeev.
  Collocation and least squares method for Navier-Stokes equations //
  Conference on Numerical Methods and Computational Mechanics 1998,
  August 24 -- 27, 1998, Miskolc, Hungary. --- Book of abstracts. ---
  Pp. 94--96. 

\end{thebibliography}

\end{document} 
